%-------------------------------------------------------------------------------
%	SECTION TITLE
%-------------------------------------------------------------------------------
\cvsection{Work Experience}


%-------------------------------------------------------------------------------
%	CONTENT
%-------------------------------------------------------------------------------
\begin{cventries}

%---------------------------------------------------------
  \cventry
    {\textbf {System Software Engineer} | \color{awesome} Tech Stack: C/C++, Rust, Secure OS, Android System Security} % Job title
    {NVIDIA} % Organization
    {Taipei, Taiwan} % Location
    {Nov. 2024 - Present} % Date(s)
    {
      \begin{cvitems} % Description(s) of tasks/responsibilities
        \item {Developed security software in the Trusted Execution Environment (TEE) for NVIDIA Tegra SoCs—spanning from user-space components to low-level device drivers—through close collaboration with hardware, platform, and application teams}
        \item {Implemented Android security HALs for cryptographic modules (Keymint, Gatekeeper) and DRM (Widevine), together with their corresponding Trusted Applications (TAs) in OP-TEE for production deployment}
        \item {Architected and built a high-performance Rust-to-C/C++ bridge for Android Keymint and OP-TEE integration using Foreign Function Interface (FFI), ensuring memory safety while maintaining optimal cryptographic performance in the TEE}
      \end{cvitems}
    }

  \vspace{12pt}
%---------------------------------------------------------
  \cventry
    {\textbf {Software Engineer Intern} | \color{awesome} Tech Stack: Modern C++, Bazel, Protobuf, Perfetto} % Job title
    {MediaTek} % Organization
    {Hsinchu, Taiwan} % Location
    {Jul. 2023 - Aug. 2023} % Date(s)
    {
      \begin{cvitems} % Description(s) of tasks/responsibilities
        \item {Enhanced the 5G modem profiler by introducing new visualization features, benefiting over 100 colleagues}
        \item {Developed a parser that serializes data into Google's Protocol Buffers format, seamlessly integrating with Google's Perfetto Trace Viewer UI, thereby enabling more effective analysis and debugging of performance issues}
      \end{cvitems}
    }

  \vspace{12pt}

%---------------------------------------------------------
  \cventry
    {\textbf{Software Engineer / Quantitative Trader} | \color{awesome} Tech Stack: Modern C++, CMake, RESTful API, WebSocket, GoogleTest, Docker} % Job title
    {A\&R Research Capital} % Organization
    {Taipei, Taiwan} % Location
    {Jan. 2022 - Sep. 2022} % Date(s)
    {
      \begin{cvitems} % Description(s) of tasks/responsibilities
        \item Led a 5-person infrastructure team to develop tools for crypto trading, utilizing a Test-Driven Development (TDD) approach
        \item initiate a project to redesigned the trading engine by converting CRTP-based static polymorphism to runtime polymorphism through proper inheritance hierarchies, resulting in improved maintainability and easier onboarding for new team members
        \item Developed a robust C++ high-frequency trading algorithm with sub-millisecond response time, yielding consistent profits 
        % \item Automated the build and testing processes by scripting, achieving a more efficient and streamlined development workflow
        % \item Reduced the system build time by over 50\% by adopting a compiler cache to avoid unnecessary recompilations
        \item Elevated test coverage from 40\% to 90\% by writing unit tests utilizing GoogleTest
      \end{cvitems}
    }

  \vspace{12pt}

%---------------------------------------------------------
\end{cventries}
